\section{Simulator Equations of Motion, Dynamics, Kinematics, and Reference Frames}
\label{sec:simulator_eom_dynamics_frames}

This section documents the ascent simulator implemented in the project, with emphasis on:
\begin{itemize}
\item the equations of motion (EOM) actually integrated,
\item the force and propulsion models,
\item the kinematic relations used to propagate position and attitude-like quantities, and
\item the reference-frame treatment (rotating-Earth vs inertial quantities).
\end{itemize}

\subsection{State Vector and Modeling Scope}

The solver integrates a reduced planar ascent model written in spherical coordinates. The base state is
\begin{equation}
\mathbf{x} = [s,\; r,\; v,\; \gamma,\; m]^T,
\end{equation}
where:
\begin{itemize}
\item $s$ is downrange distance along Earth surface [m],
\item $r$ is geocentric radius [m],
\item $v$ is vehicle speed magnitude [m/s],
\item $\gamma$ is flight-path angle (FPA) [rad],
\item $m$ is vehicle mass [kg].
\end{itemize}

When Earth rotation is enabled, latitude is additionally propagated:
\begin{equation}
\mathbf{x} = [s,\; r,\; v,\; \gamma,\; m,\; \phi]^T,
\end{equation}
and when heading tracking is enabled, heading is also propagated:
\begin{equation}
\mathbf{x} = [s,\; r,\; v,\; \gamma,\; m,\; \phi,\; \psi]^T,
\end{equation}
with $\phi$ geocentric latitude [rad] and $\psi$ local heading/azimuth [rad].

\subsection{Force and Environment Models}

\subsubsection{Gravity}

Gravity is central and inverse-square:
\begin{equation}
 g(r) = \frac{\mu_E}{r^2},
\end{equation}
with Earth gravitational parameter $\mu_E$.

\subsubsection{Atmosphere and Aerodynamics}

The atmospheric density is modeled as an exponential law:
\begin{equation}
 \rho(h) = \rho_0 \exp\!\left(-\frac{h}{H}\right),
 \qquad h = r - R_E,
\end{equation}
where $R_E$ is Earth radius, $\rho_0$ sea-level density, and $H$ scale height.

Dynamic pressure is
\begin{equation}
 q = \frac{1}{2}\rho v^2.
\end{equation}

Drag is
\begin{equation}
 F_D = q C_D A.
\end{equation}

A lift model exists,
\begin{equation}
 F_L = q C_L A,
\end{equation}
but in the current ascent propagation lift is set to zero in the dynamics call ($F_L=0$).

\subsubsection{Propulsion and Mass Flow}

Thrust magnitude $F_T$ and specific impulse $I_{sp}$ are stage-dependent and event-driven (MECO, stage separation timing, second-stage ignition, second-stage cutoff). Mass depletion follows
\begin{equation}
 \dot{m} = -\frac{F_T}{I_{sp} g_0}.
\end{equation}

Thrust direction is controlled through angle of attack (or steering command) $\alpha$, generated by the selected guidance law (gravity turn, simple polynomial, linear tangent, bilinear tangent, or Apollo-like guidance).

\subsection{Base Equations of Motion (Non-Rotating Form)}

The core differential model (before optional rotating-frame pseudo-force terms) is:
\begin{align}
 \dot{s} &= \frac{R_E}{r} v\cos\gamma, \\
 \dot{r} &= v\sin\gamma, \\
 \dot{v} &= \frac{F_T}{m}\cos\alpha - \frac{F_D}{m} - g(r)\sin\gamma, \\
 \dot{\gamma} &= \frac{1}{v}\left(\frac{F_T}{m}\sin\alpha + \frac{F_L}{m} - \left[g(r)-\frac{v^2}{r}\right]\cos\gamma\right), \\
 \dot{m} &= -\frac{F_T}{I_{sp} g_0}.
\end{align}

Implementation notes:
\begin{itemize}
\item A small-$v$ safeguard sets $\dot{\gamma}=0$ when $v$ is near zero to avoid singularity.
\item During the initial vertical hold before pitchover starts, $\dot{\gamma}$ is explicitly forced to zero.
\end{itemize}

\subsection{Kinematics and Geometry}

\subsubsection{Downrange to Latitude Mapping}

Instead of integrating latitude from a fixed local relation alone, the simulator reconstructs latitude from great-circle geometry to avoid drift and preserve bounds:
\begin{equation}
 \sigma = \frac{s}{R_E},
\end{equation}
\begin{equation}
 \sin\phi = \sin\phi_0\cos\sigma + \cos\phi_0\sin\sigma\cos\beta_0,
\end{equation}
where $\phi_0$ is launch latitude and $\beta_0$ is inertial launch azimuth.

Then
\begin{equation}
 \phi = \arcsin(\sin\phi).
\end{equation}

Latitude rate is propagated from the same geometry:
\begin{equation}
 \dot{\phi} = \frac{d\phi}{d\sigma}\frac{\dot{s}}{R_E}.
\end{equation}

\subsubsection{Heading Propagation (Optional)}

When heading is in the state, its rate is computed from spherical-geometric consistency:
\begin{equation}
 \dot{\psi} = \tan\phi\,\tan\psi\,\dot{\phi}.
\end{equation}

\subsection{Reference Frames and Earth Rotation Treatment}

\subsubsection{Frames Used}

The simulator uses two practical frames:
\begin{itemize}
\item a local rotating-Earth frame for most ascent integration quantities ($v$, $\gamma$, heading decomposition),
\item an inertial frame (ECI-equivalent) for orbital-energy/orbital-element calculations and target-orbit checks.
\end{itemize}

This mixed strategy keeps guidance and atmosphere coupling simple while preserving correct orbital mechanics metrics.

\subsubsection{Launch Azimuth with Earth Rotation}

Given launch latitude $\phi$ and target inclination $i$, the geometric inertial azimuth $\beta_{in}$ is obtained from
\begin{equation}
 \sin\beta_{in} = \frac{\cos i}{\cos\phi}.
\end{equation}

A corrected rotating-frame azimuth is then computed by accounting for eastward surface speed $v_{rot}=\Omega_E R_E\cos\phi$ and circular target speed $v_{orb}=\sqrt{\mu_E/r_{target}}$:
\begin{align}
 v_N &= v_{orb}\cos\beta_{in}, \\
 v_{E,rot} &= v_{orb}\sin\beta_{in} - v_{rot}, \\
 \beta_{corr} &= \operatorname{atan2}(v_{E,rot}, v_N).
\end{align}

Depending on settings, the active ascent azimuth is either $\beta_{corr}$ (default corrected mode) or $\beta_{in}$ (geometric mode).

\subsubsection{ECEF-like to ECI Velocity Conversion}

For orbital checks and final element computation, local rotating-frame velocity components are transformed to inertial by adding Earth rotational eastward speed:
\begin{align}
 v_h &= v\cos\gamma, \\
 v_N &= v_h\cos\psi, \qquad v_E = v_h\sin\psi, \\
 v_{E,ECI} &= v_E + \Omega_E r\cos\phi, \\
 v_{h,ECI} &= \sqrt{v_{E,ECI}^2 + v_N^2}, \\
 v_{ECI} &= \sqrt{v_{h,ECI}^2 + (v\sin\gamma)^2}, \\
 \gamma_{ECI} &= \operatorname{atan2}(v\sin\gamma,\; v_{h,ECI}).
\end{align}

\subsubsection{Optional Rotating-Frame Pseudo-Forces in EOM}

When enabled, Coriolis and centrifugal accelerations are projected into the 2D ascent dynamics. In local east-north-up (ENU) components:
\begin{align}
 a^{cor}_E &= -2\Omega_E\cos\phi\,v_{up} + 2\Omega_E\sin\phi\,v_N, \\
 a^{cor}_N &= -2\Omega_E\sin\phi\,v_E, \\
 a^{cor}_{up} &= 2\Omega_E\cos\phi\,v_E,
\end{align}
\begin{align}
 a^{cen}_E &= 0, \\
 a^{cen}_N &= -\Omega_E^2 r\sin\phi\cos\phi, \\
 a^{cen}_{up} &= \Omega_E^2 r\cos^2\phi.
\end{align}

Total pseudo-acceleration is
\begin{equation}
 \mathbf{a}_{ps} = \mathbf{a}^{cor}+\mathbf{a}^{cen}.
\end{equation}

Projecting onto along-heading horizontal direction:
\begin{equation}
 a_{h,\parallel} = a_E\sin\psi + a_N\cos\psi.
\end{equation}

Then the additive terms applied to the base EOM are:
\begin{align}
 \Delta\dot{v} &= a_{h,\parallel}\cos\gamma + a_{up}\sin\gamma, \\
 \Delta\dot{\gamma} &= \frac{-a_{h,\parallel}\sin\gamma + a_{up}\cos\gamma}{v}.
\end{align}

The integrated equations become
\begin{equation}
 \dot{v}_{tot}=\dot{v}_{base}+\Delta\dot{v}, \qquad
 \dot{\gamma}_{tot}=\dot{\gamma}_{base}+\Delta\dot{\gamma}.
\end{equation}

\subsection{Orbital Mechanics Quantities Used by the Simulator}

After conversion to inertial $(v_{ECI},\gamma_{ECI})$, orbital elements are computed with two-body relations:
\begin{align}
 a &= \frac{\mu_E r}{2\mu_E-r v_{ECI}^2}, \\
 e &= \sqrt{1-\frac{(r v_{ECI}\cos\gamma_{ECI})^2}{\mu_E a}}, \\
 r_{apo} &= a(1+e), \qquad r_{peri}=a(1-e), \\
 T &= 2\pi\sqrt{\frac{a^3}{\mu_E}}.
\end{align}

These are used for event checks (e.g., target-apogee matching in single-burn optimization) and end-of-run orbital reporting.

\subsection{Guidance Coupling with the Dynamics}

The guidance subsystem provides steering command $\alpha(t)$ (and in Apollo mode optionally thrust magnitude command). Dynamics remains the same structure, but with time-varying control inputs:
\begin{itemize}
\item before kick start: near-vertical ascent,
\item initial pitchover: time-shaped kick profile,
\item post-atmosphere-exit: selected explicit guidance law updates $\alpha$ at a configured update rate,
\item coast intervals: $F_T=0$ so only gravity, drag (if relevant), and optional pseudo-forces remain.
\end{itemize}

Thus, the simulator is a controlled nonlinear ODE with event-triggered mode switching (staging, ignition/cutoff, guidance transitions, and trajectory interrupts).

\subsection{Assumptions and Limitations}

Key assumptions to keep in mind when interpreting results:
\begin{itemize}
\item Spherical Earth with central gravity (no $J_2$ or higher harmonics).
\item Exponential atmosphere and fixed aerodynamic coefficients.
\item Reduced 2D ascent representation with optional latitude/heading augmentation.
\item Earth rotation handled through azimuth correction, inertial conversion, and optional pseudo-force injection, rather than full 3D rigid-body propagation.
\item Lift is available in formulation but currently neglected in integrated force balance.
\end{itemize}

Within these assumptions, the model is consistent for ascent guidance comparison, staging logic evaluation, and near-orbital insertion studies.
